% ============================================
% DOC-10
% Gerado de molambudos.md — build_miolo.py
% ============================================
\section*{\textbf{DOC-10} --- Registro de Eletroconvulsoterapia}
\markboth{DOC-10 --- Registro de Eletroconvulsoterapia}{}

\noindent \textit{37 sessões sem anestesia}

\noindent \textbf{Hospital Colônia de Barbacena --- Ala de Terapias Agressivas}
\textbf{Paciente:} Joaquim Antônio Correia --- Registro 1.260
\textbf{Período:} 1937 a 1952

\vspace{0.5em}\noindent\rule{\textwidth}{0.5pt}\vspace{0.5em}

\subsection*{SUMÁRIO}

\noindent Sessões registradas: \textbf{37}
Média anual: \textbf{2.46}
Equipamento: ECT Konvulsator Modelo IV (Alemanha, 1932)
Anestesia: \textbf{Não utilizada} (nenhuma das 37 sessões)
Relaxante muscular: \textbf{Não utilizado} (nenhuma das 37 sessões)
Sedação prévia: \textbf{Não} (paciente classificado como \textquotedblleft{}refratário à sedação verbal\textquotedblright{})

\vspace{0.5em}\noindent\rule{\textwidth}{0.5pt}\vspace{0.5em}

\subsection*{REGISTRO DETALHADO (parcial)}

\begin{moltablesix}
\hline
\# & Data & Vol\-ta\-gem & Duração & Efeitos Imediatos & Efeitos Tardios \\
\hline
1 & 12/mar/37 & 110V & 2.5s & Convulsão tônico-clônica, espuma, perda de consciência por 12min & Laceração de língua, incontinência urinária, confusão por 6h \\
2 & 19/mar/37 & 110V & 3.0s & Convulsão, fratura de vértebra T7 (não tratada) & Dor lombar crônica, relatada pelo paciente por 40+ anos \\
3 & 26/mar/37 & 130V & 3.0s & Convulsão prolongada (4min), cianose & Amnésia retrógrada de 48h \\
4 & 2/abr/37 & 130V & 3.5s & Convulsão, mordedura grave de língua & Dificuldade de alimentação por 10 dias \\
5 & 9/abr/37 & 130V & 3.5s & Convulsão, deslocamento de mandíbula & Redução manual sem anestesia \\
--- & --- & --- & --- & --- & --- \\
18 & 15/jun/40 & 150V & 4.0s & Convulsão, parada respiratória (reanimada) & Paciente relata \textquotedblleft{}flutuar para fora do corpo\textquotedblright{} \\
--- & --- & --- & --- & --- & --- \\
27 & 3/nov/45 & 150V & 4.0s & Convulsão, fratura de fêmur direito & Claudicação permanente \\
--- & --- & --- & --- & --- & --- \\
37 & 18/fev/52 & 170V & 5.0s & Convulsão máxima, parada cardíaca (reanimado) & Paciente inconsciente por 3 dias. Ao acordar: \textquotedblleft{}Joaquim não está aqui\textquotedblright{} \\
\hline
\end{moltablesix}

\vspace{0.5em}\noindent\rule{\textwidth}{0.5pt}\vspace{0.5em}

\subsection*{OBSERVAÇÕES CLÍNICAS (transcrição)}

\noindent \textbf{Dr. Nelson Fontes, 1937:}
\textit{\textquotedblleft{}Paciente não responde à terapia. Mantém delírio inalterado. Aumentar voltagem para 130V.\textquotedblright{}}

\noindent \textbf{Dr. Amílcar Barbosa, 1940:}
\textit{\textquotedblleft{}Após 18 sessões, paciente passou a relatar 'experiências extracorpóreas' durante o choque. Considero isto um avanço --- a mente está se desligando do corpo. Continuar.\textquotedblright{}}

\noindent \textbf{Dr. Amílcar Barbosa, 1945:}
\textit{\textquotedblleft{}Fratura de fêmur durante sessão. Paciente não sentiu a fratura --- relatou apenas que 'ouviu o osso quebrar mas não doeu'. Sugiro aumentar voltagem. Tolerância do paciente é excepcional.\textquotedblright{}}

\noindent \textbf{Dr. Heitor Oliveira, 1979 (nota marginal, 30 anos depois):}
\textit{\textquotedblleft{}Esta ficha é um documento de tortura. 37 sessões sem anestesia. Fraturas não tratadas. Sem consentimento. E o paciente sobreviveu a tudo --- lúcido o suficiente para escrever um diário de 62 anos. Isso não é refratariedade. Isso é resiliência. Ou possessão. Ou as duas coisas.\textquotedblright{}}

\vspace{0.5em}\noindent\rule{\textwidth}{0.5pt}\vspace{0.5em}

\subsection*{NOTA DO ARQUIVISTA (2026)}

\noindent O equipamento Konvulsator Modelo IV foi projetado para aplicação com anestesia e relaxante muscular. No Colônia, foi usado sem nenhum dos dois. A voltagem máxima recomendada pelo fabricante era 130V. O paciente 1.260 recebeu até 170V na última sessão.

\noindent A sobrevivência do paciente sem dano cerebral aparente é \textquotedblleft{}clinicamente inexplicável\textquotedblright{} (Dr. Torres, 2023).

\noindent O paciente 1.260 declarou, em seu diário (1952): \textit{\textquotedblleft{}O choque não afeta a criatura. O choque queima pedaços de mim. A criatura come o que sobra.\textquotedblright{}}

\noindent Sugere-se que a \textquotedblleft{}inexplicável\textquotedblright{} tolerância do paciente ao tratamento pode estar relacionada ao conteúdo de seu delírio --- especificamente, à crença em uma entidade interna que absorvia o trauma.

\noindent Ou, alternativamente, que o delírio não era delírio.

\noindent O arquivo não toma posição sobre esta última hipótese.

\vspace{0.5em}
\noindent\textit{\textbf{↪ Rotas:} \rota{MEM-17} \rotasep \rota{DOC-02} \rotasep \rota{EPI-01}}